\chapter{Related Work}\label{ch:related_work}
    % Generated by Claude
    
    % =================================================================
    % RULE: This chapter argues, it does not list.
    % Every subsection groups work BY THEME, not by author/chronology.
    % Every subsection ends with a sentence naming a limitation that
    % motivates what YOU do. If a paragraph has no such sentence,
    % delete the paragraph or fix it.
    % =================================================================

    This chapter situates the present work within existing research.
    Section~\ref{sec:rw-theme1}--\ref{sec:rw-theme3} group prior work by
    approach rather than by publication date, and Section~\ref{sec:rw-gap}
    states explicitly which gap this thesis/report addresses.
    
    % -----------------------------------------------------------------
    \section{Theme One}
    \label{sec:rw-theme1}
    
    % Group 2-4 papers here that share an approach or assumption.
    % BAD:  "Smith et al. (2020) proposed X. Jones et al. (2021) proposed Y."
    % GOOD: "Several approaches address \ac{AI}-based classification by
    %        assuming labeled training data is available
    %        \cite{smith2020,jones2021}. This assumption holds in
    %        controlled settings but breaks down when [your actual
    %        problem context], because [concrete reason]."
    
    Placeholder: describe the shared approach of this cluster of work
    \cite{example2024}, then state its limitation relative to your problem.
    
    % -----------------------------------------------------------------
    \section{Theme Two}
    \label{sec:rw-theme2}
    
    Placeholder: second thematic cluster. Do not restate Theme One's
    limitation here — every theme needs a distinct angle, or merge the
    themes.
    
    % -----------------------------------------------------------------
    \section{Theme Three}
    \label{sec:rw-theme3}
    
    Placeholder: third thematic cluster, if you actually have one. Do not
    pad this to hit three sections — two solid themes beat three thin ones.
    
    % -----------------------------------------------------------------
    \section{Comparison of Approaches}
    \label{sec:rw-comparison}
    
    Table~\ref{tab:rw-comparison} summarizes the reviewed approaches along
    the dimensions most relevant to this work. Choose dimensions that
    actually differentiate the work — if every row has the same value in
    a column, delete that column.
    
    \begin{table}[H]
        \centering
        \caption{Comparison of related approaches.}
        \label{tab:rw-comparison}
        \begin{tabularx}{\textwidth}{l X X X}
            \toprule
            Approach & Method & Assumption & Limitation \\
            \midrule
            \cite{example2024} & Placeholder & Placeholder & Placeholder \\
            \bottomrule
        \end{tabularx}
    \end{table}

    % -----------------------------------------------------------------
    \section{Positioning of This Work}
    \label{sec:rw-gap}
    
    % MANDATORY. This subsection cannot be a summary of the above.
    % It must state, in plain declarative sentences:
    %   1. What gap exists across ALL the themes above combined.
    %   2. Why closing that gap matters (not "would be interesting" —
    %      concretely, for whom, and what breaks without it).
    %   3. One sentence on how this work closes it.
    % If you cannot write these three sentences, you have not found
    % a gap yet — go back and re-read the literature, don't fake it here.
    
    Placeholder: state the gap. Do not write "no prior work has
    considered X" unless you have actually verified this by searching,
    not by hoping.